\section{Mục tiêu}
\begin{frame}{Dự án có thể tái lập}
\begin{itemize}
\item Nêu một câu hỏi và kết quả đo được.
\item Tách bằng chứng khỏi giả định.
\item Giữ dữ liệu và công cụ cùng báo cáo.
\end{itemize}
\end{frame}
\section{Ví dụ}
\begin{frame}{Kết quả tổng quan}
\centering\includegraphics[width=.65\linewidth,height=.55\textheight,keepaspectratio]{figures/results.png}
\par Bốn lượt chạy; các số đo tăng tuyến tính.
\end{frame}
\begin{frame}{Luồng xử lý}
\centering\begin{tikzpicture}
\node[draw,rounded corners] (a) at (0,0) {Đầu vào};
\node[draw,rounded corners] (b) at (3,0) {Kiểm tra};
\node[draw,rounded corners] (c) at (6,0) {Đầu ra};
\draw[->,thick] (a)--(b);\draw[->,thick] (b)--(c);\end{tikzpicture}
\end{frame}
\begin{frame}{Kết luận}
Ví dụ minh họa cấu trúc, không phải khẳng định hiệu năng. Hãy thay dữ liệu mẫu, sau đó báo cáo sai số và giới hạn.
\[T(n)=4n\]
\end{frame}
